\documentclass[]{beamer}
\usetheme{KUL}
\usepackage{multirow}
\usepackage{multicol}
\usepackage{tikz}
\usepackage{ulem}
\usepackage{siunitx}
\newcommand\itemS{\item[\textbf{\S}]}
\definecolor{darkgreen}{rgb}{0,0.598,0.199}
\usepackage{times} % set font on times new roman
\usepackage{eurosym} % package for Euro sign
\usepackage{lineno}   % package for line numbering
\usepackage{hyperref} % this is for url links
\usepackage{subcaption}  % this package enables one to put several figures next to each other
\usepackage{textcomp}
\usepackage{setspace}
\usepackage{gensymb}
\usepackage{hyperref}
\usepackage[backend=biber,style=apa]{biblatex}
\addbibresource{ref.bib}

%----------------------------------
% Fill in the essential Information
%----------------------------------

\title[Main title]{Knots and fractals}
\subtitle{\ldots a subtitle}
\author[A.\ Depicker]{Arthur Depicker} % between [] is short name, between {} is long name
\date{3/4/2025} % Here you can also just type something, e.g. October 10, 2017
\institute[KU Leuven]{Faculty of Science\\ Department of Earth and Environmental Sciences}

%----------------------------------
% ACTUAL PRESENTATION STARTS HERE
%----------------------------------

\begin{document}

% TITEL PAGE	
	{
		\setbeamertemplate{headline}{} %define local, empty header for title page
		\setbeamertemplate{footline}{} %define local, empty footer for title page
		\maketitle
	}
	\addtocounter{framenumber}{-1} % We don't count the title page

% FRAME 1 
\section{Knot}
\begin{frame}{What is a knot}
\begin{figure}
\centering
\includegraphics[width=7cm]{Pictures/images.jpg}
\includegraphics[width=3cm]{Pictures/chineseknot.png}
\caption{Knots in daily life (\cite{Knotsdailylife} and \cite{Chineseknot})}
    \label{fig:enter-label}
\end{figure}
\end{frame}

% FRAME 2
\begin{frame}{What is a knot}
    \begin{itemize}
        \item Link
    \end{itemize}
\begin{figure}
\centering
\includegraphics[width = 4.9cm]{Pictures/audi.jpg}
\includegraphics[width = 4cm]{Pictures/borromean.png}
    \caption{Links in logos(\cite{audi} and \cite{christiantrequetra})}
    \label{fig:enter-label}
\end{figure}


\end{frame}

% FRAME 3 
\begin{frame}{What is a knot}
\begin{itemize}
    \item To study knots rigorously, we need to connect the two ends of the rope.
\end{itemize}
\begin{kulblock}{Definition}
    In mathematics, a knot is an embedding of the circle $\mathbf{S}^1$ into 3-dimensional Euclidean space, $\mathbf{R}^3$.
\end{kulblock}

\end{frame}

% FRAME 4
\begin{frame}{Motivation}
\begin{itemize}
    \item Peter Guthrie Tait and Lord Kelvin (19th century): During the research on composition of atoms, they conjectured that atoms are made out of vortex ring of ether.
\end{itemize}
\begin{figure}
    \centering
\includegraphics[width=5cm]{Pictures/simple knots.png}
    \caption{Simple knots(\cite{simpleknots})}
    \label{fig:enter-label}
\end{figure}

\end{frame}

% FRAME 5
\begin{frame}{Motivation}
\begin{itemize}
    \item The structure of protein and DNA
\end{itemize}
\begin{figure}
    \centering
\includegraphics[width=7cm]{Pictures/DNA.png}
    \caption{DNA (\cite{DNA})}
    \label{fig:enter-label}
\end{figure}
\centering
\includegraphics[width=7cm]{Pictures/DNA.png}
\end{frame}

% FRAME 6
\begin{frame}{Motivation}
\begin{itemize}
    \item The knot structure of chemical compound and material
\end{itemize}
\begin{figure}
    \centering
\includegraphics[width = 4.5cm]{Pictures/chemical.png}
\includegraphics[width=4.5cm]{Pictures/material.jpg}
    \caption{Tightest synthetic compound (\cite{compound} and \cite{farbric})}
    \label{fig:enter-label}
\end{figure}

\end{frame}

% FRAME 7
\begin{frame}{Diagram}
\begin{itemize}
    \item How to represent a knot? Regular projection.
\end{itemize}
\begin{kulblock}{Definition}
    A knot projection is called a regular projection if no three points on the knot project to the same point, and no vertex
projects to the same point as any other point on the knot.
\end{kulblock}
    
\end{frame}

% FRAME 8
\begin{frame}{Diagram}

\begin{itemize}
    \item A knot diagram is the regular projection of a knot to the plane with broken
lines indicating where one part of the knot undercrosses the other part.
\end{itemize}
\begin{figure}
    \centering
\includegraphics[width=3cm]{Pictures/trefoil.png}
\includegraphics[width=3cm]{Pictures/31.png}
    \caption{Trefoil diagram (\cite{trefoil})}
    \label{fig:enter-label}
\end{figure}
\end{frame}

% FRAME 9
\begin{frame}{Equivalence of knots}
\begin{itemize}
    \item Are they equivalent?
\end{itemize}
\begin{figure}
   \centering
    \includegraphics[width = 4cm]{Pictures/Simple.png}
    \includegraphics[width = 4cm]{Pictures/complex.png}
    \caption{Knot unfolding (\cite{knotsunfolding})}
    \label{fig:enter-label}
\end{figure}

\end{frame}


% FRAME 10
\begin{frame}{Reidemeister move}
\begin{itemize}
    \item Equivalence of knot?
\end{itemize}
\begin{figure}
   \centering
\includegraphics[width=7cm, height = 2cm]{Pictures/equivalence.png} 
    \caption{Equivalence? (\cite{Equivalence})}
    \label{fig:enter-label}
\end{figure}


\begin{itemize}
    \item Formal definition: two knots $\mathbf{K}_1$ and $\mathbf{K}_2$ are equivalent if there is an orientation-preserving homeomorphism $h:\mathbb{R}^3 \rightarrow \mathbb{R}^3$ with $h(\mathbf{K}_1) = h(\mathbf{K}_2)$.
\end{itemize}

\begin{itemize}
    \item (1927)Reidemeister's Theorem(\cite{Reidemeister}): If two knots are equivalent, their diagrams are related by a sequence of Reidemeister moves.
\end{itemize}
\end{frame}

% FRAME 11
\begin{frame}{Reidemeister move}
\begin{itemize}
    \item A Reidemeister move is an operation that can be performed on the diagram
of a knot whithout altering the corresponding knot.
\end{itemize}
\begin{figure}
    \centering
\includegraphics[width=7cm]{Pictures/r.png}
    \caption{Equivalence? (\cite{reidemeister})}
    \label{fig:enter-label}
\end{figure}

\end{frame}

% FRAME 12
\begin{frame}{Invariant polynomial}
    \begin{itemize}
        \item 1923(James Waddell Alexander)(\cite{Alexander}): Alexander polynomial $\triangledown$. Reformulated by John Conway in 1969.
\begin{figure}
   \centering
\includegraphics[width =7cm, height = 2.5cm]{Pictures/L.png}
    \caption{3 variations (\cite{L+})}
    \label{fig:enter-label}
\end{figure}
        $$\triangledown(O) = 1$$
        $$\triangledown(L_+) - \triangledown(L_-) = z\triangledown(L_0)$$
    \end{itemize}
\end{frame}

% FRAME 13
\begin{frame}{Invariant polynomial}
\begin{itemize}
    \item 1984(Vaughan Jones)(\cite{Jones}): Jones polynomial $V$. Jones won the field's medal for this discovery in 1990.
    $$V(L_0) = 1$$
    $$t^{-1}V(L_+) - tV(L_-) = (t^{\frac{1}{2}} - t^{-\frac{1}{2}})V(L_0)$$
    \item 1985(multiple authors): HOMFLY-PT polynomial $P$: A generalization of the above invariants.
    $$P(L_0) = 1$$
    $$tP(L_+) + t^{-1}P(L_-) + mP(L_0) = 0$$
\end{itemize}
\end{frame}


% FRAME 14
\begin{frame}{Fractals}
\begin{figure}
    \centering
\includegraphics[width = 8cm]{Pictures/fractalcover.jpg}
    \caption{Fractal (\cite{fractalcover})}
    \label{fig:enter-label}
\end{figure}

\end{frame}

% FRAME 15
\begin{frame}{Motivation}
\begin{itemize}
    \item Coastline
    \item Lungs
\begin{figure}
    \centering
    \includegraphics[width = 5cm]{Pictures/coastline.jpg}
    \includegraphics[height = 3.4cm, width  = 4cm]{Pictures/lungs.jpg}
    \caption{Coastline and lung (\cite{COastline} and \cite{Lung})}
    \label{fig:enter-label}
\end{figure}
    
\end{itemize}
\end{frame}



% FRAME 16
\begin{frame}{Motivation}
    \begin{itemize}
        \item 1872: Karl Weierstrass constructed a function that is continuous but nowhere differentiable:
        $$f(x) = \sum\limits_{n=0}^{\infty}a^ncos(b^n\pi x)$$
    \end{itemize}
    \begin{itemize}
        \item 1904: Helge von Koch devised a geometric method to construct pathological curves.
    \end{itemize}
    \begin{itemize}
        \item 1915: Wacław Sierpiński introduced Sierpiński triangle.
    \end{itemize}
    \begin{itemize}
        \item 1918: Gaston Julia & Pierre Fatou studied iteration of complex functions.
    \end{itemize}
\end{frame}

% FRAME 17
\begin{frame}[allowframebreaks]{References}
\nocite{*}
\printbibliography
\end{frame}


\end{document}